% GENERATED FILE. Edit canonical lessons, then run npm run generate:books.
% canonical-source-hash: fnv1a64:21c522228a1b5d53
% canonical-lessons: TE-C467-antastu, TE-C467-nelamaliga, TE-C467-prahari, TE-C467-visalamaina, TE-C467-nirmanusyamaina

\chapter{Buildings and Spaces}
\label{ch:te-buildings-and-spaces}
\hlchaptermodality{467}

\begin{chapteropening}
\textbf{By the end of this chapter:} \emph{I can say a storey, a basement, a compound wall, spacious and deserted.}

Complete the last lesson of chapter 467: I can say a storey, a basement, a compound wall, spacious and deserted.
\end{chapteropening}

\section[\texorpdfstring{\te{అంతస్తు} (antastu)}{అంతస్తు (antastu)}]{\te{అంతస్తు} (antastu) — a storey; a floor}

\label{lesson:TE-C467-antastu}



\begin{quote}
Before the new one: say the Telugu for a street lamp, then the Telugu for surroundings.
\end{quote}

\subsection*{You'll want to know: \te{అంతస్తు}}
\textbf{\te{అంతస్తు}} — \emph{antastu} — \textquotedblleft{}a storey; a floor\textquotedblright{}.

\begin{quote}
\te{మా} \te{ఇల్లు} \te{మూడో} \te{అంతస్తులో} \te{ఉంది}.
\end{quote}

\begin{quote}
\emph{mā illu mūḍō antastulō undi.}
\end{quote}

\begin{quote}
Our flat is on the third floor.
\end{quote}

\te{కింది} \te{అంతస్తు}, the ground floor; \te{పై} \te{అంతస్తు}, the upper floor. A five-storey building is \te{ఐదు} \te{అంతస్తుల} \te{భవనం}.

\begin{cousinweb}
\textbf{\te{అంతస్తు}} (\emph{antastu}) is a Telugu word for a storey whose origin is not settled.
\end{cousinweb}

\begin{grammarlens}[title={What you've built}]
One more word for getting around.
\end{grammarlens}

\subsection*{Your turn}
\begin{itemize}
\raggedright
  \item \emph{Say it:} \emph{antastu}
  \item \emph{Say it:} \emph{antastu}, once more
  \item \emph{Say it:} \emph{mā illu mūḍō antastulō undi.} — Our flat is on the third floor.
  \item \emph{Recall:} say the Telugu for a street lamp, then the Telugu for surroundings, then say \emph{antastu} again
\end{itemize}

\begin{tcolorbox}[breakable,colback=teal!4,colframe=teal!35!black,title={Before you move on}]
What does \te{అంతస్తు} mean? (\textquotedblleft{}A storey; a floor\textquotedblright{}.) Say it once more.
\end{tcolorbox}
\section[\texorpdfstring{\te{నేలమాళిగ} (nēlamāḷiga)}{నేలమాళిగ (nēlamāḷiga)}]{\te{నేలమాళిగ} (nēlamāḷiga) — a basement; a cellar}

\label{lesson:TE-C467-nelamaliga}



\begin{quote}
Before the new one: say the Telugu for surroundings, then the Telugu for a storey.
\end{quote}

\subsection*{You'll want to know: \te{నేలమాళిగ}}
\textbf{\te{నేలమాళిగ}} — \emph{nēlamāḷiga} — \textquotedblleft{}a basement; a cellar\textquotedblright{}.

\begin{quote}
\te{ఆ} \te{భవనం} \te{నేలమాళిగలో} \te{పార్కింగు} \te{ఉంది}.
\end{quote}

\begin{quote}
\emph{ā bhavanaṁ nēlamāḷigalō pārkiṅgu undi.}
\end{quote}

\begin{quote}
There is parking in the basement of that building.
\end{quote}

\te{నేల}, the ground, + \te{మాళిగ}, a house: a room below the ground. Old palaces kept treasure in one.

\begin{cousinweb}
\textbf{\te{నేలమాళిగ}} (\emph{nēlamāḷiga}) is native Dravidian: \te{నేల}, \textquotedblleft{}the earth, the ground\textquotedblright{} (DEDR 2913), and \te{మాళిగ}, \textquotedblleft{}a house, a storied house\textquotedblright{} (DEDR 4796).
\end{cousinweb}

\begin{grammarlens}[title={What you've built}]
One more word for getting around.
\end{grammarlens}

\subsection*{Your turn}
\begin{itemize}
\raggedright
  \item \emph{Say it:} \emph{nēlamāḷiga}
  \item \emph{Say it:} \emph{nēlamāḷiga}, once more
  \item \emph{Say it:} \emph{ā bhavanaṁ nēlamāḷigalō pārkiṅgu undi.} — There is parking in the basement of that building.
  \item \emph{Recall:} say the Telugu for surroundings, then the Telugu for a storey, then say \emph{nēlamāḷiga} again
\end{itemize}

\begin{tcolorbox}[breakable,colback=teal!4,colframe=teal!35!black,title={Before you move on}]
What does \te{నేలమాళిగ} mean? (\textquotedblleft{}A basement; a cellar\textquotedblright{}.) Say it once more.
\end{tcolorbox}
\section[\texorpdfstring{\te{ప్రహరీ} (praharī)}{ప్రహరీ (praharī)}]{\te{ప్రహరీ} (praharī) — a compound wall}

\label{lesson:TE-C467-prahari}



\begin{quote}
Before the new one: say the Telugu for a storey, then the Telugu for a basement.
\end{quote}

\subsection*{You'll want to know: \te{ప్రహరీ}}
\textbf{\te{ప్రహరీ}} — \emph{praharī} — \textquotedblleft{}a compound wall\textquotedblright{}.

\begin{quote}
\te{బడి} \te{చుట్టూ} \te{ఎత్తైన} \te{ప్రహరీ} \te{గోడ} \te{ఉంది}.
\end{quote}

\begin{quote}
\emph{baḍi cuṭṭū ettaina praharī gōḍa undi.}
\end{quote}

\begin{quote}
There is a high compound wall round the school.
\end{quote}

Mostly \te{ప్రహరీ} \te{గోడ}: the wall round a house, a school or an office.

\begin{cousinweb}
\textbf{\te{ప్రహరీ}} (\emph{praharī}) is a Telugu word for a boundary wall whose origin is not settled.
\end{cousinweb}

\begin{grammarlens}[title={What you've built}]
One more word for getting around.
\end{grammarlens}

\subsection*{Your turn}
\begin{itemize}
\raggedright
  \item \emph{Say it:} \emph{praharī}
  \item \emph{Say it:} \emph{praharī}, once more
  \item \emph{Say it:} \emph{baḍi cuṭṭū ettaina praharī gōḍa undi.} — There is a high compound wall round the school.
  \item \emph{Recall:} say the Telugu for a storey, then the Telugu for a basement, then say \emph{praharī} again
\end{itemize}

\begin{tcolorbox}[breakable,colback=teal!4,colframe=teal!35!black,title={Before you move on}]
What does \te{ప్రహరీ} mean? (\textquotedblleft{}A compound wall\textquotedblright{}.) Say it once more.
\end{tcolorbox}
\section[\texorpdfstring{\te{విశాలమైన} (viśālamaina)}{విశాలమైన (viśālamaina)}]{\te{విశాలమైన} (viśālamaina) — spacious; broad}

\label{lesson:TE-C467-visalamaina}



\begin{quote}
Before the new one: say the Telugu for a basement, then the Telugu for a compound wall.
\end{quote}

\subsection*{You'll want to know: \te{విశాలమైన}}
\textbf{\te{విశాలమైన}} — \emph{viśālamaina} — \textquotedblleft{}spacious; broad\textquotedblright{}.

\begin{quote}
\te{హోటలులో} \te{మాకు} \te{విశాలమైన} \te{గది} \te{ఇచ్చారు}.
\end{quote}

\begin{quote}
\emph{hōṭalulō māku viśālamaina gadi iccāru.}
\end{quote}

\begin{quote}
At the hotel they gave us a spacious room.
\end{quote}

It stands before its noun: \te{విశాలమైన} \te{రోడ్డు}, a broad road. \te{విశాలంగా} is the adverb.

\begin{cousinweb}
\textbf{\te{విశాలమైన}} (\emph{viśālamaina}) is built from Sanskrit \emph{viśāla}, \textquotedblleft{}spacious, broad\textquotedblright{} (CDIAL 11936), with the Telugu adjective ending -\te{మైన}.
\end{cousinweb}

\begin{grammarlens}[title={What you've built}]
One more word for getting around.
\end{grammarlens}

\subsection*{Your turn}
\begin{itemize}
\raggedright
  \item \emph{Say it:} \emph{viśālamaina}
  \item \emph{Say it:} \emph{viśālamaina}, once more
  \item \emph{Say it:} \emph{hōṭalulō māku viśālamaina gadi iccāru.} — At the hotel they gave us a spacious room.
  \item \emph{Recall:} say the Telugu for a basement, then the Telugu for a compound wall, then say \emph{viśālamaina} again
\end{itemize}

\begin{tcolorbox}[breakable,colback=teal!4,colframe=teal!35!black,title={Before you move on}]
What does \te{విశాలమైన} mean? (\textquotedblleft{}Spacious; broad\textquotedblright{}.) Say it once more.
\end{tcolorbox}
\section[\texorpdfstring{\te{నిర్మానుష్యమైన} (nirmānuṣyamaina)}{నిర్మానుష్యమైన (nirmānuṣyamaina)}]{\te{నిర్మానుష్యమైన} (nirmānuṣyamaina) — deserted; empty of people}

\label{lesson:TE-C467-nirmanusyamaina}



\begin{quote}
Before the new one: say the Telugu for a compound wall, then the Telugu for spacious.
\end{quote}

\subsection*{You'll want to know: \te{నిర్మానుష్యమైన}}
\textbf{\te{నిర్మానుష్యమైన}} — \emph{nirmānuṣyamaina} — \textquotedblleft{}deserted; empty of people\textquotedblright{}.

\begin{quote}
\te{అర్ధరాత్రి} \te{నేను} \te{నిర్మానుష్యమైన} \te{వీధిలో} \te{నడిచాను}.
\end{quote}

\begin{quote}
\emph{ardharātri nēnu nirmānuṣyamaina vīdhilō naḍicānu.}
\end{quote}

\begin{quote}
At midnight I walked down a deserted street.
\end{quote}

\te{నిర్}, without, + \te{మానుష్య}, people. With \te{ఉండు} it takes -\te{ంగా}: \te{నిర్మానుష్యంగా} \te{ఉంది}.

\begin{cousinweb}
\textbf{\te{నిర్మానుష్యమైన}} (\emph{nirmānuṣyamaina}) is built from Sanskrit \emph{nir-}, \textquotedblleft{}without\textquotedblright{}, and \emph{mānuṣya}, \textquotedblleft{}human, the state of men\textquotedblright{} (Monier-Williams), with the Telugu adjective ending -\te{మైన}.
\end{cousinweb}

\begin{grammarlens}[title={What you've built}]
One more word for getting around.
\end{grammarlens}

\subsection*{Your turn}
\begin{itemize}
\raggedright
  \item \emph{Say it:} \emph{nirmānuṣyamaina}
  \item \emph{Say it:} \emph{nirmānuṣyamaina}, once more
  \item \emph{Say it:} \emph{ardharātri nēnu nirmānuṣyamaina vīdhilō naḍicānu.} — At midnight I walked down a deserted street.
  \item \emph{Recall:} say the Telugu for a compound wall, then the Telugu for spacious, then say \emph{nirmānuṣyamaina} again
\end{itemize}

\begin{tcolorbox}[breakable,colback=teal!4,colframe=teal!35!black,title={Before you move on}]
What does \te{నిర్మానుష్యమైన} mean? (\textquotedblleft{}Deserted; empty of people\textquotedblright{}.) Say it once more.
\end{tcolorbox}
